\chapter{Algebra homologiczna: kompleksy i (ko)homologie}
\label{ch:algebra-homologiczna}
\index{homologia@Homologia}

Grupy, moduły i tensorowanie zebraliśmy w
Rozdziale~\ref{ch:algebra-abstrakcyjna}.
Teraz badamy ciągi odwzorowań między modułami.
W kompleksie łańcuchowym elementy spełniające
$\partial x=0$ są cyklami; część z nich ma postać
$x=\partial y$ i jest brzegami. Iloraz cykli
przez brzegi tworzy homologię. W kohomologii
strzałki biegną w przeciwną stronę.
Na razie prowadzimy rachunek algebraiczny;
później zastosujemy go do przestrzeni.

\section{Ciągi modułów i dokładność}
\label{sec:dokladnosc-modulow}

Mając kolejne homomorfizmy
\[
 \cdots\longrightarrow M_{i-1}
 \xrightarrow{f_{i-1}}M_i
 \xrightarrow{f_i}M_{i+1}
 \longrightarrow\cdots,
\]
pytamy, które elementy $M_i$ przychodzą z lewej
i które znikają po wysłaniu w prawo.
Samo słowo \emph{dokładny} odnosi się do
ciągu \emph{w danym miejscu}, a nie do
pojedynczego modułu.

\begin{definicja}[Dokładność ciągu]
\index{dokladnosc ciagu@Dokładność ciągu}
\label{def:ciag-dokladny}
Ciąg jest \emph{dokładny w $M_i$}, jeśli
\[
 \operatorname{im}f_{i-1}=\ker f_i.
\]
Jest \emph{dokładny}, jeśli tak jest
w każdym jego miejscu. Równość zawiera
dwa żądania: $f_i f_{i-1}=0$ oraz to,
że każdy element zabity przez $f_i$
już pochodzi z $M_{i-1}$.
\emph{Długi ciąg dokładny} ma wiele
następujących po sobie takich miejsc;
może rozciągać się w obie strony
bez wyznaczonego początku.
\end{definicja}

Ponieważ jedyne odwzorowanie $0\to M$ ma obraz $0$,
warunek dokładności $0\to A\xrightarrow{i}B$
w $A$ mówi $\ker i=0$, czyli $i$ jest iniekcją.
Analogicznie dokładność
$B\xrightarrow{p}C\to0$ w $C$ mówi
$\operatorname{im}p=C$, czyli $p$ jest surjekcją.

\begin{definicja}[Krótki ciąg dokładny]
\index{krotki ciag dokladny@Krótki ciąg dokładny}
\label{def:krotki-ciag}
Ciąg
\begin{equation}\label{eq:krotki-ciag}
 0\longrightarrow A\xrightarrow{i}B
   \xrightarrow{p}C\longrightarrow0
\end{equation}
jest \emph{krótki i dokładny}, gdy $i$ jest
iniekcją, $p$ jest surjekcją oraz
$\operatorname{im}i=\ker p$.
Wtedy $C\cong B/i(A)$ przez
Twierdzenie~\ref{thm:modul-izomorfizm}
zastosowane do $p$. Mówimy też,
że $B$ jest \emph{rozszerzeniem} $C$ przez $A$.
\end{definicja}

\begin{przyklad}[Reszty z dzielenia]
\label{prz:ses-mod-m}
Dla $m\geq2$ mamy krótki ciąg dokładny
\[
 0\longrightarrow\Z\xrightarrow{\ \times m\ }\Z
   \xrightarrow{\ \pi\ }\Z/m\Z\longrightarrow0,
 \qquad \pi(k)=[k].
\]
Pierwsze odwzorowanie jest iniektywne: $mk=0$
w $\Z$ pociąga $k=0$. Odwzorowanie $\pi$ jest
surjektywne, bo każda klasa ma reprezentanta.
Jej jądro to dokładnie wielokrotności $m$,
czyli obraz pierwszego odwzorowania.
\end{przyklad}

Krótki ciąg może, ale nie musi, rozkładać się
na dwie niezależne części. Nie należy mylić
wyboru dowolnego przedstawiciela każdej klasy
z \emph{liniowym} wyborem przedstawicieli.

\begin{twierdzenie}[Kiedy krótki ciąg się rozszczepia?]
\label{thm:ses-rozszczepienie}
Dla~\eqref{eq:krotki-ciag} równoważne są:
\begin{enumerate}[leftmargin=*]
\item istnieje homomorfizm $s:C\to B$
      z $p s=\id_C$ (przekrój $p$);
\item istnieje homomorfizm $r:B\to A$
      z $r i=\id_A$ (retrakcja $i$);
\item istnieje izomorfizm $B\cong A\oplus C$
      zgodny z odwzorowaniami $i(a)=(a,0)$
      i $p(a,c)=c$.
\end{enumerate}
Wtedy ciąg nazywamy \emph{rozszczepionym}.
\end{twierdzenie}
\begin{proof}
Załóżmy najpierw, że mamy $s$.
Każde $b\in B$ można rozłożyć jako
\[
 b=(b-s(p(b)))+s(p(b)).
\]
Pierwszy składnik leży w $\ker p=\operatorname{im}i$.
Ponieważ $i$ jest iniekcją, istnieje dokładnie
jeden $a\in A$ taki, że $i(a)=b-s(p(b))$.
Odwzorowanie
$\Phi:A\oplus C\to B$, $\Phi(a,c)=i(a)+s(c)$
jest więc surjektywne. Jeżeli $\Phi(a,c)=0$,
to po zastosowaniu $p$ otrzymujemy $c=0$,
a wtedy $i(a)=0$, czyli $a=0$.
To dowodzi (1)$\Rightarrow$(3).

Z izomorfizmu $\Phi:A\oplus C\to B$ w (3) i rzutu
$\operatorname{pr}_A:A\oplus C\to A$ bierzemy
$r=\operatorname{pr}_A\circ\Phi^{-1}$. Dostajemy
$r:B\to A$ z $ri=\id_A$. To daje (3)$\Rightarrow$(2).
Wreszcie, mając $r$, dla $c\in C$ wybierzmy
jakikolwiek $b$ z $p(b)=c$ i połóżmy
$s(c)=b-i(r(b))$. Jeśli zmienimy $b$ na
$b+i(a)$, nowa wartość wyniesie
\[
 b+i(a)-i(r(b+i(a)))
 =b+i(a)-i(r(b)+a)=b-i(r(b)),
\]
ponieważ $r(i(a))=a$ i $r$ jest liniowa.
Tak otrzymane odwzorowanie $s$ nie zależy od
wybranego $b$ i jest liniowe: dla $c,c'$
możemy jako podniesienie $c+c'$ wziąć sumę
podniesień $c,c'$, a dla $tc$ podniesienie $tb$.
Ponadto $p(s(c))=p(b)=c$.
Otrzymaliśmy (2)$\Rightarrow$(1).
\end{proof}

\begin{uwaga}[Dwa kontrastujące przypadki]
Jeśli $R=\mathbb K$ jest ciałem, każdy krótki
ciąg przestrzeni liniowych rozszczepia się:
wybierzmy bazę $C$, podnieśmy jej wektory
do $B$ przez surjekcję $p$ i rozszerzmy
liniowo do odwzorowania $s$. Nie zależy nam na
kanoniczności tego wyboru.

Ciąg z Przykładu~\ref{prz:ses-mod-m} dla
$m\geq2$ \emph{nie} rozszczepia się jako
ciąg $\Z$-modułów. Gdyby przekrój $s$
istniał, dla $z=s([1])\in\Z$ byłoby
$mz=s(m[1])=0$, więc $z=0$.
Ale wtedy $\pi(s([1]))=[0]\ne[1]$.
Brak dzielenia przez $m$ uniemożliwia tu
liniowy wybór reprezentantów.
\end{uwaga}
\section{Kompleksy łańcuchowe i homologia algebraiczna}
\label{sec:kompleks-homologia}

W dokładnym ciągu każdy element zabity przez
następne odwzorowanie pochodzi z poprzedniego miejsca.
Do zdefiniowania homologii wystarczy najpierw
słabszy warunek: obraz poprzedniego odwzorowania ma
\emph{mieścić się} w jądrze następnego.
Homologia mierzy brak równości.

\begin{definicja}[Kompleks łańcuchowy]
\index{kompleks lancuchowy@Kompleks łańcuchowy}
\label{def:kompleks-lancuchowy}
\emph{Kompleks łańcuchowy} $C_\bullet$ nad $R$
to moduły $C_n$, $n\in\Z$, i homomorfizmy
$\partial_n:C_n\to C_{n-1}$ takie, że
\begin{equation}\label{eq:d-kwadrat-zero}
 \partial_{n-1}\circ \partial_n=0\qquad\text{dla każdego }n.
\end{equation}
Odwzorowanie $\partial_n$ nazywamy \emph{różniczką} albo
\emph{brzegiem}. Dla kompleksów łańcuchowych
będziemy używać symbolu $\partial$, a symbol $d$
pozostawimy różniczkom kołańcuchowym.
Gdy $C_n=0$ dla $n<0$,
mówimy o kompleksie w stopniach nieujemnych.
\end{definicja}

Warunek~\eqref{eq:d-kwadrat-zero} należy czytać
dosłownie: brzeg tego, co samo jest brzegiem,
wynosi zero. Nie zakładamy, że wszystkie
elementy o zerowym brzegu są brzegami.

\begin{definicja}[Cykle, brzegi i grupy homologii]
\index{cykle brzegi i grupy homologii@Cykle, brzegi i grupy homologii}
\label{def:homologia-algebraiczna}
W stopniu $n$ określamy podmoduły
\[
 Z_n(C):=\ker(\partial_n:C_n\to C_{n-1}),\qquad
 B_n(C):=\operatorname{im}(\partial_{n+1}:C_{n+1}\to C_n).
\]
Elementy $Z_n(C)$ to \emph{cykle}, a elementy
$B_n(C)$ to \emph{brzegi}. Ponieważ
$\partial_n\partial_{n+1}=0$, mamy $B_n(C)\subseteq Z_n(C)$,
więc można zdefiniować \emph{homologię stopnia $n$}
\begin{equation}\label{eq:homologia-def}
 H_n(C):=Z_n(C)/B_n(C)
 =\frac{\ker \partial_n}{\operatorname{im}\partial_{n+1}}.
\end{equation}
Klasę cyklu $z$ oznaczamy $[z]$.
\end{definicja}

Dwa cykle $z,z'$ opisują tę samą klasę dokładnie
wtedy, gdy $z-z'=\partial_{n+1}w$ dla pewnego
$w\in C_{n+1}$. To jest algebraiczne
znaczenie słowa \emph{homologiczne}.
Klasa $[z]$ znika dokładnie wtedy, gdy
$z$ jest brzegiem. Z Definicji~\ref{def:ciag-dokladny}
otrzymujemy zatem
\[
 H_n(C)=0
 \quad\Longleftrightarrow\quad
 \operatorname{im}\partial_{n+1}=\ker \partial_n
 \quad\Longleftrightarrow\quad
 C_\bullet\text{ jest dokładny w }C_n.
\]
Gdy znika homologia we wszystkich stopniach,
kompleks nazywamy \emph{acyklicznym}.
Nie znaczy to, że same moduły $C_n$ są zerowe.

\begin{przyklad}[Cztery rachunki]
\label{prz:homologia-rachunki}
\begin{enumerate}[leftmargin=*]
\item W kompleksie
      $0\to\Z\xrightarrow{\ \times2\ }\Z\to0$,
      gdzie niezerowe moduły są w stopniach
      $1$ i $0$, mamy $Z_1=0$, $Z_0=\Z$,
      $B_0=2\Z$. Stąd
      $H_1=0$ i $H_0=\Z/2\Z$.
\item Jeżeli na tych samych dwóch modułach
      postawimy różniczkę zerową, każdy element
      jest cyklem, a jedynym brzegiem jest $0$.
      Otrzymujemy $H_1\cong H_0\cong\Z$.
\item Dla
      $0\to R\xrightarrow{\ \id\ }R\to0$
      odwzorowanie w stopniu $1$ ma zerowe jądro,
      a w stopniu $0$ jego obraz jest całym $R$.
      Wszystkie homologie są zerowe, choć
      dwa moduły kompleksu są niezerowe.
\item Torsja może powstać nawet z modułów
      wolnych. Weźmy
      \[
       0\longrightarrow\Z
       \xrightarrow{\ \partial_2\ }\Z^2
       \xrightarrow{\ \partial_1\ }\Z
       \longrightarrow0,\qquad
       \partial_2(t)=(0,2t),\quad \partial_1(x,y)=x.
      \]
      Mamy $\partial_1\partial_2=0$.
      Cykle w stopniu $1$ to $(0,y)$,
      a brzegi to $(0,2t)$, więc
      $H_1\cong\Z/2\Z$. Odwzorowanie $\partial_1$ jest
      surjektywne, więc $H_0=0$; odwzorowanie
      $\partial_2$ jest iniektywne, więc $H_2=0$.
\end{enumerate}
\end{przyklad}

\begin{definicja}[Odwzorowanie łańcuchowe]
\index{odwzorowanie lancuchowe@Odwzorowanie łańcuchowe}
\label{def:mapa-lancuchowa}
Odwzorowanie $f:C_\bullet\to D_\bullet$ to rodzina
homomorfizmów $f_n:C_n\to D_n$ zgodnych
z różniczkami:
\begin{equation}\label{eq:mapa-lancuchowa}
 \partial^D_n f_n=f_{n-1}\partial^C_n.
\end{equation}
Oznacza to, że najpierw zastosować $f$,
a potem wziąć brzeg, daje ten sam wynik,
co zrobić to w odwrotnej kolejności.
\end{definicja}

\begin{stwierdzenie}[Odwzorowanie indukowane na homologii]
\label{prop:homologia-funktor}
Odwzorowanie łańcuchowe $f$ wyznacza homomorfizmy
$H_n(f):H_n(C)\to H_n(D)$,
$H_n(f)[z]=[f_n(z)]$. Zachodzą
$H_n(\id_C)=\id_{H_n(C)}$ oraz
$H_n(gf)=H_n(g)H_n(f)$.
\end{stwierdzenie}
\begin{proof}
Jeżeli $\partial^C_n z=0$, to z~\eqref{eq:mapa-lancuchowa}
$\partial^D_n f_nz=f_{n-1}\partial^C_nz=0$, czyli
$f_nz$ jest cyklem. Jeżeli zmienimy
$z$ na $z+\partial^C_{n+1}w$, jego obraz zmieni się o
\[
 f_n(\partial^C_{n+1}w)
 =\partial^D_{n+1}(f_{n+1}w),
\]
czyli o brzeg. Odwzorowanie na klasach jest więc
dobrze określone i liniowe.
Równości dla identyczności i złożenia
wynikają przez zastosowanie wzorów
do dowolnej klasy $[z]$.
\end{proof}

\section{Skąd bierze się długi ciąg dokładny homologii?}
\label{sec:dlugi-ciag-homologii}

Niech $A_\bullet,B_\bullet,C_\bullet$ będą
kompleksami, a $i:A_\bullet\to B_\bullet$,
$p:B_\bullet\to C_\bullet$ odwzorowaniami łańcuchowymi.
Mówimy, że
\[
 0\longrightarrow A_\bullet
 \xrightarrow{i}B_\bullet
 \xrightarrow{p}C_\bullet
 \longrightarrow0
\]
jest \emph{krótkim ciągiem dokładnym kompleksów},
jeśli dla \emph{każdego} $n$ ciąg modułów
\begin{equation}\label{eq:ses-kompleks}
 0\longrightarrow A_n\xrightarrow{i_n}B_n
 \xrightarrow{p_n}C_n\longrightarrow0
\end{equation}
jest dokładny. Stopniowa dokładność nie
mówi jeszcze, że ciąg powstałych homologii
jest krótki i dokładny: cykl podniesiony
z $C_n$ może przestać być cyklem w $B_n$.
Jego brzeg mierzy właśnie tę przeszkodę.

\begin{figure}[!htbp]
\centering
\begin{tikzpicture}[>=Stealth,line cap=round,
  every node/.style={font=\small}]
 \node[draw=manifoldblue!75!black,rounded corners,
  fill=manifoldblue!7,minimum width=2.25cm,
  minimum height=.9cm] (c) at (-5.25,0) {$c\in Z_n(C)$};
 \node[draw=manifoldgreen!75!black,rounded corners,
  fill=manifoldgreen!8,minimum width=2.05cm,
  minimum height=.9cm] (b) at (-1.75,0) {$b\in B_n$};
 \node[draw=manifoldred!75!black,rounded corners,
  fill=manifoldred!8,minimum width=2.15cm,
  minimum height=.9cm] (db) at (1.75,0) {$\partial_Bb\in B_{n-1}$};
 \node[draw=manifoldgold!85!black,rounded corners,
  fill=manifoldgold!10,minimum width=2.4cm,
  minimum height=.9cm] (a) at (5.25,0) {$a\in Z_{n-1}(A)$};
 \draw[->,thick,manifoldblue] (c) --
  node[midway,above,font=\footnotesize] {wybór $b$} (b);
 \draw[->,thick,manifoldred] (b) --
  node[midway,above,font=\footnotesize] {$\partial_B$} (db);
 \draw[->,thick,manifoldgold!85!black] (db) --
  node[midway,above,font=\footnotesize] {$i_{n-1}^{-1}$} (a);
 \node[font=\footnotesize,manifoldgreen!75!black]
  at (-1.75,-.9) {$p_n(b)=c$};
 \node[font=\footnotesize,manifoldred!85!black]
  at (1.75,-.9) {$p_{n-1}(\partial_Bb)=\partial_Cc=0$};
 \node[font=\footnotesize,manifoldgold!85!black]
  at (5.25,-.9) {$[a]=\delta_n[c]$};
\end{tikzpicture}
\caption{Odwzorowanie łączące: podnosimy cykl $c$, bierzemy brzeg
podniesienia $b$, po czym odczytujemy go w $A_{n-1}$.
Odwrócenie $i_{n-1}$ dotyczy tylko jego obrazu
$\ker p_{n-1}$; całego odwzorowania $i_{n-1}$ nie odwracamy.}
\label{fig:homologia-delta}
\end{figure}

\begin{twierdzenie}[Długi ciąg dokładny]
\label{thm:les-homologia}
Z krótkiego ciągu dokładnego kompleksów
otrzymujemy naturalny długi ciąg dokładny
\begin{equation}\label{eq:les-homologia}
 \cdots\longrightarrow H_n(A)
 \xrightarrow{i_*}H_n(B)
 \xrightarrow{p_*}H_n(C)
 \xrightarrow{\delta_n}H_{n-1}(A)
 \xrightarrow{i_*}H_{n-1}(B)
 \longrightarrow\cdots.
\end{equation}
Odwzorowanie \emph{łączące} $\delta_n$ ma następujący
konkretny wzór. Dla cyklu $c\in C_n$ wybieramy
podniesienie $b\in B_n$ z $p_n(b)=c$.
Wtedy $\partial_B b=i_{n-1}(a)$ dla dokładnie
jednego $a\in A_{n-1}$, a
\begin{equation}\label{eq:delta-wzor}
 \delta_n[c]=[a]\in H_{n-1}(A).
\end{equation}
Jeżeli kompleksy są zerowe w ujemnych stopniach,
końcowy fragment ma postać
$H_0(A)\to H_0(B)\to H_0(C)\to0$.
\end{twierdzenie}
\begin{proof}
\emph{Konstrukcja i cykl.} Surjektywność $p_n$
daje podniesienie $b$. Ponieważ $p$ jest odwzorowaniem łańcuchowym,
\[
 p_{n-1}(\partial_Bb)=\partial_C(p_nb)=\partial_Cc=0.
\]
Z dokładności~\eqref{eq:ses-kompleks} w
stopniu $n-1$ mamy
$\partial_Bb\in\ker p_{n-1}=\operatorname{im}i_{n-1}$,
więc istnieje $a$ z $i_{n-1}a=\partial_Bb$.
Iniektywność $i_{n-1}$ daje jego jednoznaczność.
Ponadto
\[
 i_{n-2}(\partial_Aa)
 =\partial_B(i_{n-1}a)=\partial_B^2b=0.
\]
Iniektywność $i_{n-2}$ wymusza $\partial_Aa=0$.
Możemy więc wziąć klasę $[a]$.

\emph{Niezależność od podniesienia.} Jeśli $b'$ też
podnosi $c$, to $b'-b=i_n(u)$ dla pewnego
$u\in A_n$. Wtedy
\[
 \partial_Bb'=\partial_Bb+\partial_Bi_nu
 =i_{n-1}(a+\partial_Au).
\]
Nowy cykl w $A$ różni się od $a$ o brzeg,
więc wyznacza tę samą klasę.

\emph{Niezależność od reprezentanta.}
Niech $c'=c+\partial_Cv$, $v\in C_{n+1}$.
Wybierzmy podniesienie $w\in B_{n+1}$ elementu $v$.
Wtedy $b+\partial_Bw$ podnosi $c'$ i ma ten sam
brzeg co $b$, bo $\partial_B^2w=0$.
Każde inne podniesienie $c'$ różni się od tego
o element $\operatorname{im}i_n$; właśnie
sprawdzona niezależność od podniesienia
kończy dowód niezależności od wyboru reprezentanta klasy $[c]$.
Suma podniesień podnosi sumę cykli, a
podniesienie $rb$ podnosi $rc$, więc $\delta_n$
jest homomorfizmem $R$-modułów.

\emph{Dokładność przy $H_n(B)$.}
Z $pi=0$ mamy $p_*i_*=0$.
Jeśli $[b]\in H_n(B)$ spełnia
$p_*[b]=0$, to dla cyklu $b$
istnieje $v\in C_{n+1}$ z
$p_nb=\partial_Cv$. Wybierzmy podniesienie
$w\in B_{n+1}$ elementu $v$.
Wtedy $b'=b-\partial_Bw$ nadal jest cyklem,
$[b']=[b]$ i
$p_nb'=p_nb-\partial_Cv=0$. Z dokładności
w stopniu $n$ mamy $b'=i_na$.
Ponieważ $\partial_Bb'=i_{n-1}\partial_Aa=0$
i $i_{n-1}$ jest iniektywna, $a$ jest
cyklem. Stąd $[b]=i_*[a]$.

\emph{Dokładność przy $H_n(C)$.}
Jeżeli $c=p_nb$ dla cyklu $b$, możemy
w definicji $\delta_n$ wybrać to właśnie
podniesienie: $\partial_Bb=0$, więc $\delta_n[c]=0$.
Odwrotnie, załóżmy $\delta_n[c]=0$
i podnieśmy cykl $c$ do $b$.
Jeśli $\partial_Bb=i_{n-1}a$, to $[a]=0$,
czyli $a=\partial_Au$ dla pewnego $u\in A_n$.
Element $b-i_nu$ jest cyklem, bo
$\partial_B(b-i_nu)=i_{n-1}(a-\partial_Au)=0$,
a $p_n(b-i_nu)=c$. Zatem
$[c]=p_*[b-i_nu]$.

\emph{Dokładność przy $H_{n-1}(A)$.}
Dla klasy $\delta_n[c]=[a]$ mamy
$i_{n-1}a=\partial_Bb$, czyli $i_*[a]=0$.
Odwrotnie, jeżeli $a\in Z_{n-1}(A)$
i $i_*[a]=0$, istnieje $b\in B_n$
takie, że $i_{n-1}a=\partial_Bb$.
Wówczas $c=p_nb$ jest cyklem, bo
$\partial_Cc=p_{n-1}\partial_Bb=p_{n-1}i_{n-1}a=0$,
a wzór~\eqref{eq:delta-wzor} daje
$\delta_n[c]=[a]$.
Przesuwając indeks $n$ o jeden,
obejmujemy tymi trzema sprawdzeniami
każde miejsce całego ciągu~\eqref{eq:les-homologia}.
W końcu, przy założonej zerowości kompleksów w ujemnych
stopniach mamy $A_{-1}=0$, a więc
$H_{-1}(A)=0$, więc $H_0(B)\to H_0(C)$
jest surjekcją na mocy właśnie wykazanej
dokładności.
\end{proof}

\begin{uwaga}[Naturalność odwzorowania łączącego]
\label{uwaga:les-naturalnosc}
Jeżeli mamy odwzorowanie między dwoma krótkimi
ciągami kompleksów, zgodne z $i,p$ i
różniczkami, to odwzorowania na homologiach
komutują także z $\delta$.
Dokładniej, są to odwzorowania łańcuchowe
$f_A:A\to A'$, $f_B:B\to B'$, $f_C:C\to C'$ spełniające
$f_Bi=i'f_A$ oraz $p'f_B=f_Cp$.
Istotnie, podniesienie $b$ cyklu $c$ przechodzi
w podniesienie obrazu $c$, jego brzeg przechodzi
w obraz $\partial_Bb$, a wstrzyknięcie $i$
pozwala odczytać obraz tego samego $a$.
Równość na reprezentantach ma postać
$\delta'_n(H_n(f_C)[c])
=H_{n-1}(f_A)(\delta_n[c])$.
\end{uwaga}

\begin{przyklad}[Niezerowe odwzorowanie łączące]
\label{prz:delta-izomorfizm}
Niech $A_\bullet$ ma jedyny niezerowy moduł
$A_0=R$, a $C_\bullet$ jedyny niezerowy
moduł $C_1=R$. Niech $B_\bullet$ ma
$B_1=B_0=R$ i różniczkę
$\partial_B=\id_R:B_1\to B_0$.
W stopniu $0$ odwzorowanie $i_0:A_0\to B_0$
jest identycznością, a w stopniu $1$
odwzorowanie $p_1:B_1\to C_1$ jest identycznością;
pozostałe odwzorowania to jedyne możliwe odwzorowania zerowe.
W każdym stopniu otrzymujemy rozszczepiony
krótki ciąg dokładny modułów, ale
$H_n(B)=0$ we wszystkich stopniach.
Cykl $c=1\in C_1$ podnosimy do $b=1\in B_1$.
Jego brzeg to $1=i_0(1)\in B_0$,
więc $\delta_1:R\to R$ jest identycznością.
Przykład pokazuje, że stopniowe rozszczepienie
nie oznacza zgodnego z różniczkami
rozszczepienia kompleksów.
\end{przyklad}
\section{Kohomologia algebraiczna i jej długi ciąg}
\label{sec:kohomologia-algebraiczna}

W homologii różniczka obniża stopień.
Kohomologia powstaje z równie prostej
konstrukcji, w której różniczka stopień
\emph{podwyższa}. Gdy elementy kompleksu
łańcuchowego są wektorami, elementy
kompleksu dualnego są funkcjonałami:
zamiast brać brzeg wektora, oceniamy
funkcjonał na brzegu.

\begin{definicja}[Kompleks kołańcuchowy i kohomologia]
\index{kompleks kolancuchowy i kohomologia@Kompleks kołańcuchowy i kohomologia}
\label{def:kohomologia-algebraiczna}
\emph{Kompleks kołańcuchowy} $K^\bullet$ to
moduły $K^n$ i homomorfizmy
$d^n:K^n\to K^{n+1}$ z
$d^{n+1}d^n=0$ dla $n\in\Z$. Jego \emph{kohomologią stopnia $n$}
jest
\begin{equation}\label{eq:kohomologia-def}
 H^n(K):=\frac{\ker(d^n:K^n\to K^{n+1})}
 {\operatorname{im}(d^{n-1}:K^{n-1}\to K^n)}.
\end{equation}
Elementy jądra nazywamy \emph{kocyklami},
a elementy obrazu \emph{kobrzegami}.
Odwzorowanie kołańcuchowe $F:K^\bullet\to L^\bullet$
składa się z $F^n:K^n\to L^n$ spełniających
$d_L^nF^n=F^{n+1}d_K^n$.
\end{definicja}

Tak jak w dowodzie
Stwierdzenia~\ref{prop:homologia-funktor},
odwzorowanie kołańcuchowe przenosi kocykle w
kocykle i kobrzegi w kobrzegi:
$d_LFz=Fdz=0$ oraz
$F(d_Ku)=d_L(Fu)$. Wobec tego definiuje
$H^n(F)[z]=[F^nz]$. Równość
$H^n(K)=0$ mówi dokładnie, że
$K^{n-1}\to K^n\to K^{n+1}$
jest dokładny w $K^n$.

\begin{przyklad}[Kohomologia dwuwyrazowego kompleksu]
Dla $K^0=K^1=\Z$, $d^0=\times2$
i wszystkich innych modułów równych $0$
otrzymujemy $H^0(K)=0$ oraz
$H^1(K)=\Z/2\Z$. W stopniu $1$ nie ma
już różniczki wychodzącej, więc
wszystkie elementy $\Z$ są kocyklami;
kobrzegami są liczby parzyste.
\end{przyklad}

\begin{stwierdzenie}[Kompleks dualny]
\label{prop:kompleks-dualny}
Dla kompleksu łańcuchowego $C_\bullet$
i $R$-modułu $G$ połóżmy
\[
 C^n(C;G):=\operatorname{Hom}_R(C_n,G),
 \qquad d^n\varphi:=\varphi\circ \partial_{n+1}.
\]
Wtedy $C^\bullet(C;G)$ jest kompleksem
kołańcuchowym. Odwzorowanie łańcuchowe
$f:C_\bullet\to D_\bullet$ daje odwzorowanie
kołańcuchowe w odwrotną stronę,
$f^*:C^\bullet(D;G)\to C^\bullet(C;G)$,
$f^*(\varphi)=\varphi\circ f_n$.
\end{stwierdzenie}
\begin{proof}
Dla $\varphi:C_n\to G$ mamy
\[
 d^{n+1}(d^n\varphi)
 =(\varphi\circ \partial_{n+1})\circ \partial_{n+2}
 =\varphi\circ(\partial_{n+1}\partial_{n+2})=0.
\]
Z kolei dla $\varphi:D_n\to G$
oraz $x\in C_{n+1}$ równość odwzorowań łańcuchowych
z~\eqref{eq:mapa-lancuchowa} daje
\[
 (d f^*\varphi)(x)
 =\varphi(f_n(\partial_C x))
 =\varphi(\partial_D(f_{n+1}x))
 =(f^*d\varphi)(x).
\]
To dowodzi zgodności z różniczką.
\end{proof}

Słowo \emph{dualny} wymaga tu ostrożności:
zastosowanie $\operatorname{Hom}_R(-,G)$
do krótkiego ciągu dokładnego modułów
nie musi dać krótkiego ciągu dokładnego.
Na przykład przy $G=\Z$ w ciągu
$0\to2\Z\hookrightarrow\Z\to\Z/2\Z\to0$
funkcjonał $2\Z\to\Z$, $2k\mapsto k$,
nie przedłuża się do homomorfizmu
$\Z\to\Z$: musiałby wysłać $1$
na liczbę $t$ z $2t=1$.
Z tego powodu nie będziemy budować
długiego ciągu kohomologii przez
bezrefleksyjne dualizowanie
Twierdzenia~\ref{thm:les-homologia}.

\begin{twierdzenie}[Długi ciąg kohomologii]
\label{thm:les-kohomologia}
Jeżeli w każdym stopniu dokładny jest
krótki ciąg \emph{kompleksów kołańcuchowych}
\[
 0\longrightarrow A^\bullet\xrightarrow{i}B^\bullet
 \xrightarrow{p}C^\bullet\longrightarrow0,
\]
to istnieje długi ciąg dokładny
\begin{equation}\label{eq:les-kohomologia}
 \cdots\longrightarrow H^n(A)
 \xrightarrow{i_*}H^n(B)
 \xrightarrow{p_*}H^n(C)
 \xrightarrow{\delta^n}H^{n+1}(A)
 \xrightarrow{i_*}H^{n+1}(B)\longrightarrow\cdots.
\end{equation}
\end{twierdzenie}
\begin{proof}
Dla kocyklu $c\in C^n$ wybierzmy podniesienie
$b\in B^n$ z $p(b)=c$.
Wtedy $p(d_Bb)=d_Cc=0$, więc dokładność
w stopniu $n+1$ i iniektywność $i$
dają dokładnie jedno $a\in A^{n+1}$
z $i(a)=d_Bb$.
Ponieważ $i(d_Aa)=d_B^2b=0$,
otrzymujemy kocykl $a$.
Definiujemy $\delta^n[c]=[a]$.
Jeśli podniesienie $b$ zastąpimy $b+i(u)$,
to $a$ zmienia się na $a+d_Au$.
Jeśli $c$ zastąpimy $c+d_Cv$, to po
wyborze podniesienia $w$ dla $v$ możemy
zastąpić $b$ przez $b+d_Bw$,
co w ogóle nie zmienia $d_Bb$.
Oba wybory prowadzą zatem do tej
samej klasy. Podniesienie sumy można wziąć
jako sumę podniesień, a $rb$ podnosi $rc$ dla $r\in R$.
Zatem $\delta^n$ jest $R$-liniowa.

Sprawdźmy każde z trzech rodzajów miejsc
w~\eqref{eq:les-kohomologia}.
Przy $H^n(B)$ mamy $p_*i_*=0$.
Jeśli $b\in B^n$ jest kocyklem
i $p_*[b]=0$, to $p(b)=d_Cv$
dla pewnego $v\in C^{n-1}$.
Podnieśmy $v$ do $w\in B^{n-1}$.
Wówczas $b-d_Bw$ jest kocyklem
z tą samą klasą co $b$, a jego
obraz przez $p$ wynosi $0$.
Równa się więc $i(a)$ dla pewnego
$a\in A^n$. Z $0=d_Bi(a)=i(d_Aa)$
i iniektywności $i$ dostajemy
$d_Aa=0$, a zatem $[b]=i_*[a]$.

Przy $H^n(C)$ inkluzja
$\operatorname{im}p_*\subseteq\ker\delta^n$
wynika z wybrania jako podniesienia $c=p(b)$
samego kocyklu $b$: jego różniczka jest $0$.
Jeśli odwrotnie $\delta^n[c]=0$,
to dla podniesienia $b$ mamy $d_Bb=i(a)$,
gdzie $a=d_Au$ dla pewnego $u\in A^n$.
Element $b-i(u)$ jest kocyklem,
którego obraz to $c$, więc
$[c]\in\operatorname{im}p_*$.

Przy $H^{n+1}(A)$ element
$i_*\delta^n[c]$ reprezentuje
$i(a)=d_Bb$ i dlatego jest zerowy.
Jeśli z drugiej strony $a\in A^{n+1}$
jest kocyklem i $i_*[a]=0$,
to $i(a)=d_Bb$ dla pewnego
$b\in B^n$. Obraz $c=p(b)$ jest
kocyklem, gdyż $d_Cc=p(d_Bb)=p(i(a))=0$,
a jego odwzorowanie łączące wynosi $[a]$.
Zmiana $n$ daje dokładność także
w pozostałych stopniach. Konstrukcja jest naturalna:
dla zgodnych odwzorowań $f_A,f_B,f_C$ między krótkimi ciągami
podniesieniem $f_C(c)$ jest $f_B(b)$, a jego różniczka to
$f_B(d_Bb)=f_B(i(a))=i'(f_A(a))$.
Stąd $\delta'^n H^n(f_C)=H^{n+1}(f_A)\delta^n$.
\end{proof}

\begin{twierdzenie}[Nad ciałem kohomologia jest dualna do homologii]
\label{thm:field-hom-dual}
Niech $C_\bullet$ będzie kompleksem
przestrzeni liniowych nad ciałem
$\mathbb K$, a
$C^n=\operatorname{Hom}_{\mathbb K}(C_n,\mathbb K)$.
Wtedy naturalne przypisanie
\begin{equation}\label{eq:field-duality}
 H^n(C^\bullet)
 \longrightarrow\operatorname{Hom}_{\mathbb K}
  (H_n(C_\bullet),\mathbb K),
 \qquad [\varphi]\longmapsto
     ([z]\longmapsto\varphi(z))
\end{equation}
jest izomorfizmem. Zakładamy, że
przestrzenie pozwalają na wybór
dopełnień podprzestrzeni
(w szczególności wystarcza skończony wymiar).
\end{twierdzenie}
\begin{proof}
Jeżeli $\varphi\in C^n$ jest kocyklem,
to $d^n\varphi=\varphi \partial_{n+1}=0$.
Zatem $\varphi$ znika na $B_n$,
więc jej wartość na $z\in Z_n$
zależy tylko od klasy $[z]\in H_n$.
Jeżeli zmienimy $\varphi$ na
$\varphi+d^{n-1}\psi
=\varphi+\psi \partial_n$, to dla cyklu $z$
mamy $\partial_nz=0$, więc otrzymujemy
ten sam funkcjonał na $H_n$.
Odwzorowanie~\eqref{eq:field-duality} jest
dobrze określone i liniowe. Jest też naturalne: dla odwzorowania
łańcuchowego $f:C\to D$, kocyklu $\varphi$ na $D$ i cyklu $z$ na $C$
mamy $(f^*\varphi)(z)=\varphi(fz)$. Zatem obie drogi, przez
$H^n(f^*)$ lub przez prekompozycję z $H_n(f)$, dają ten sam funkcjonał.

Aby pokazać surjektywność, weźmy
$\ell:H_n=Z_n/B_n\to\mathbb K$.
Złożenie $Z_n\to H_n\xrightarrow{\ell}\mathbb K$
jest liniowe i znika na $B_n$.
Wybieramy dopełnienie $Z_n$ w $C_n$
i przedłużamy funkcjonał liniowo
na $C_n$, na przykład jako zero
na tym dopełnieniu. Otrzymany
$\varphi$ spełnia
$\varphi \partial_{n+1}=0$, bo
$\operatorname{im}\partial_{n+1}=B_n$;
jest kocyklem wysyłanym na $\ell$.

Jeżeli klasa $[\varphi]$ trafia na zero,
to $\varphi$ znika na całym $Z_n$.
Możemy więc zdefiniować funkcjonał
$\eta$ na $B_{n-1}=\operatorname{im}\partial_n$
przez $\eta(\partial_nx)=\varphi(x)$.
Jest dobrze określony: z
$\partial_nx=\partial_nx'$ wynika $x-x'\in Z_n$,
a tam $\varphi$ znika.
Wybieramy dopełnienie $B_{n-1}$
w $C_{n-1}$ i przedłużamy $\eta$
do $\psi:C_{n-1}\to\mathbb K$.
Wtedy $\varphi=\psi \partial_n=d^{n-1}\psi$
jest kobrzegiem. Jądro odwzorowania
w~\eqref{eq:field-duality} jest więc zerowe.
\end{proof}

\begin{uwaga}[Dlaczego ważne są współczynniki?]
Nad $\Z$ odpowiednik~\eqref{eq:field-duality}
może być fałszywy. Dla kompleksu
$C_1=C_0=\Z$ z $\partial_1=\times2$ mamy
$H_0(C)=\Z/2$ i $H_1(C)=0$,
podczas gdy dla dualnego kompleksu
$H^1(C^\bullet)=\Z/2$.
Grupa $\operatorname{Hom}_{\Z}
(H_1(C),\Z)$ jest zerowa.
Ta dodatkowa klasa kohomologii
wykrywa informację torsyjną z
poprzedniego stopnia. Nie utożsamiamy
więc homologii z kohomologią bez
sprawdzenia założeń o współczynnikach.
\end{uwaga}
\section{Homotopia łańcuchowa i uproszczenia kompleksów}
\label{sec:homotopia-lancuchowa}

Jedna homologia może mieć bardzo różne opisy
przez generatory i różniczki. Poniższe kryterium
mówi, kiedy dwa odwzorowania kompleksów dają
\emph{dokładnie to samo} odwzorowanie na
homologiach. Jest to algebraiczny mechanizm
usuwania par generatorów, które wzajemnie
się kasują.

\begin{definicja}[Homotopia łańcuchowa]
\index{homotopia lancuchowa@Homotopia łańcuchowa}
\label{def:homotopia-lancuchowa}
Dwa odwzorowania łańcuchowe $f,g:C_\bullet\to D_\bullet$
są \emph{łańcuchowo homotopijne}, jeśli istnieją
homomorfizmy $R$-modułów $h_n:C_n\to D_{n+1}$ takie, że
\begin{equation}\label{eq:chain-homotopy}
 f_n-g_n=\partial^D_{n+1}h_n+h_{n-1}\partial^C_n
 \qquad\text{dla każdego }n.
\end{equation}
Kompleks $C_\bullet$ jest \emph{ściągalny},
jeśli $\id_C$ jest łańcuchowo homotopijna
z odwzorowaniem zerowym.
\end{definicja}

Wzór~\eqref{eq:chain-homotopy} przypomina
zasadę całkowania przez części: różnica
odwzorowań składa się z dwóch operacji, w których
$h$ i $\partial$ występują w przeciwnych
kolejnościach. Na cyklach drugi składnik
znika, a pierwszy jest brzegiem.

\begin{twierdzenie}[Homotopia nie zmienia homologii]
\label{thm:homotopia-homologia}
Jeśli $f$ i $g$ są łańcuchowo homotopijne,
to $H_n(f)=H_n(g)$ dla każdego $n$.
W szczególności kompleks ściągalny jest
acykliczny. Jeśli istnieją odwzorowania łańcuchowe
$f:C_\bullet\to D_\bullet$ i
$g:D_\bullet\to C_\bullet$ z
$gf\simeq\id_C$ i $fg\simeq\id_D$,
to $H_n(f)$ i $H_n(g)$ są wzajemnymi
izomorfizmami.
\end{twierdzenie}
\begin{proof}
Dla $z\in Z_n(C)$ równanie~\eqref{eq:chain-homotopy}
redukuje się do
\[
 f_n(z)-g_n(z)=\partial^D_{n+1}h_n(z),
\]
bo $\partial^C_nz=0$. Różnica obrazów jest
brzegiem, więc $[f_n(z)]=[g_n(z)]$.
Gdy $f=\id_C$, $g=0$, otrzymujemy
$[z]=[0]$ dla każdej klasy, czyli
$H_n(C)=0$.
Na koniec, dzięki
Stwierdzeniu~\ref{prop:homologia-funktor},
\[
 H_n(g)H_n(f)=H_n(gf)
 =H_n(\id_C)=\id,\qquad
 H_n(f)H_n(g)=\id.
\]
Stąd oba odwzorowania na homologiach są
odwrotne do siebie.
\end{proof}

\begin{przyklad}[Najprostsze ściągnięcie]
Dla kompleksu
$0\to R\xrightarrow{\id}R\to0$
z modułami w stopniach $1$ i $0$
połóżmy $h_0:R=C_0\to C_1=R$
równe identyczności, a pozostałe
$h_n=0$. W stopniu $0$ mamy
$\partial_1h_0=\id_{C_0}$, a w stopniu
$1$ mamy $h_0\partial_1=\id_{C_1}$.
Wzór~\eqref{eq:chain-homotopy}
daje $\id_C=\partial h+h\partial$.
\end{przyklad}

\begin{twierdzenie}[Algebraiczne skasowanie pary generatorów]
\label{thm:kasowanie-pary}
Niech $C_\bullet$ będzie kompleksem
wolnych $R$-modułów z ustalonymi
bazami, a $n\geq1$. Załóżmy, że
przy generatorze $f\in C_{n-1}$
we współrzędnych $\partial_n(e)$ dla pewnego
generatora $e\in C_n$ stoi skalar
odwracalny $u\in R$.
Wtedy przez zmiany baz można rozłożyć
$C_\bullet$ na sumę prostą
podkompleksu i dwuwyrazowego kompleksu
\[
 0\longrightarrow R e
 \xrightarrow{\ \partial_n\ }R\,\partial_n(e)
 \longrightarrow0,
\]
w którym $\partial_n$ jest izomorfizmem.
Usunięcie tej pary nie zmienia homologii.
\end{twierdzenie}
\begin{proof}
Zapiszmy
$C_n=R e\oplus M$ oraz
$C_{n-1}=R f\oplus N$, gdzie
$M,N$ są rozpięte przez pozostałe
wybrane wektory bazowe.
Z założenia
$\partial_n(e)=u f+v$ z $v\in N$.
Ponieważ $u$ jest odwracalny,
$f':=\partial_n(e)$ może zastąpić $f$
w bazie $C_{n-1}$:
$C_{n-1}=R f'\oplus N$.
Rzeczywiście, $f=u^{-1}(f'-v)$
pozwala przejść w obie strony
między rozkładami.

Dla $x\in M$ zapiszmy
$\partial_n(x)=a(x)f+w(x)$,
gdzie $a(x)\in R$ i $w(x)\in N$.
Zmieńmy też bazę w stopniu $n$,
zastępując każdy bazowy $x\in M$
wektorem $x':=x-u^{-1}a(x)e$.
Wtedy
\[
 \partial_n(x')
 =a(x)f+w(x)-u^{-1}a(x)(uf+v)
 =w(x)-u^{-1}a(x)v\in N.
\]
Nowe wektory $x'$ razem z $e$
tworzą bazę $C_n$: możemy odzyskać
$x=x'+u^{-1}a(x)e$.
Oznaczmy ich rozpiętość przez $M'$.
Mamy zatem $\partial_n(Re)=Rf'$
i $\partial_n(M')\subseteq N$.

Równość $\partial_{n-1}(f')=\partial_{n-1}\partial_n(e)=0$
mówi, że z $Rf'$ nie wychodzi
różniczka. Z kolei dla $y\in C_{n+1}$
zapiszmy $\partial_{n+1}(y)=t e+x'$,
$x'\in M'$. Po zastosowaniu $\partial_n$
dostajemy
$0=\partial_n\partial_{n+1}(y)=t f'+\partial_n(x')$.
Pierwszy składnik należy do $Rf'$,
drugi do $N$; rozkład jest prosty,
więc $t=0$. Żadna różniczka
z wyższego stopnia nie trafia do $Re$.
Pozostałe moduły wraz z $M'$ i $N$
tworzą więc podkompleks, a
$Re\xrightarrow{\cong}Rf'$
jest jego prostym składnikiem
w całym kompleksie.

Na tym dwuwyrazowym składniku
homotopia $h_{n-1}(f')=e$,
zerowa w innych stopniach,
spełnia $\id=\partial h+h\partial$.
Jest on ściągalny. Przedłużmy tę homotopię przez zero na
pozostałym podkompleksie $D$. Dla rzutu $p:C\to D$ i włączenia
$j:D\to C$ otrzymujemy $pj=\id_D$ oraz
$\id_C-jp=\partial h+h\partial$. Rzut i włączenie
wyznaczają odwrotne izomorfizmy
na homologiach na mocy
Twierdzenia~\ref{thm:homotopia-homologia}.
\end{proof}

Nad $\Z$ jednostkami są tylko $1$ i $-1$.
Współczynnik $2$ w różniczce
$\Z\xrightarrow{\times2}\Z$ nie pozwala
wykonać takiego skasowania: pozostała
klasa $\Z/2$ z
Przykładu~\ref{prz:homologia-rachunki}
jest rzeczywistą przeszkodą.

\begin{stwierdzenie}[Acykliczny kompleks przestrzeni liniowych]
\label{prop:acykliczny-sciagalny}
Jeżeli $C_\bullet$ jest kompleksem
przestrzeni liniowych nad ciałem
$\mathbb K$ i $H_n(C)=0$ dla wszystkich
$n$, to $C_\bullet$ jest ściągalny.
Wystarczy, aby każda podprzestrzeń
$Z_n(C)\subseteq C_n$ miała dopełnienie;
warunek ten zachodzi zawsze przy
skończonych wymiarach.
\end{stwierdzenie}
\begin{proof}
Acykliczność daje $Z_n=B_n$.
Wybierzmy dopełnienie $S_n$ tak, aby
$C_n=Z_n\oplus S_n$. Ograniczenie
\[
 \partial_n|_{S_n}:S_n\longrightarrow Z_{n-1}
\]
jest iniektywne: element jego jądra
należy do $S_n\cap Z_n=\{0\}$.
Jest też surjektywne, bo
$\partial_n(C_n)=B_{n-1}=Z_{n-1}$
i $\partial_n$ znika na $Z_n$.
Ma więc odwrotność.
Każde $x\in C_n$ zapiszmy
jednoznacznie jako $x=z+s$
z $z\in Z_n$, $s\in S_n$, i połóżmy
\[
 h_n(z+s)
 :=(\partial_{n+1}|_{S_{n+1}})^{-1}(z)
 \in S_{n+1}.
\]
Wówczas $\partial_{n+1}h_n(z+s)=z$.
Z drugiej strony $\partial_n(z+s)=\partial_n(s)$,
więc
$h_{n-1}\partial_n(z+s)
=(\partial_n|_{S_n})^{-1}(\partial_ns)=s$.
Zatem $(\partial h+h\partial)(x)=z+s=x$
w każdym stopniu.
\end{proof}

\begin{przyklad}[Acykliczny, lecz nieściągalny nad $\Z$]
\label{prz:acykliczny-niesciagalny}
Potraktujmy dokładny ciąg z
Przykładu~\ref{prz:ses-mod-m} dla $m=2$
jako kompleks w stopniach $2,1,0$:
\[
 0\longrightarrow\Z\xrightarrow{\ \times2\ }\Z
 \xrightarrow{\ \pi\ }\Z/2\Z\longrightarrow0.
\]
Jego wszystkie homologie są zerowe,
bo $\times2$ jest iniekcją,
$\ker\pi=2\Z$, a $\pi$ jest
surjekcją. Gdyby kompleks był
ściągalny, to w stopniu $0$
wzór $\id=\partial h+h\partial$ dałby
$\id_{\Z/2}=\pi h_0$.
Odwzorowanie $h_0:\Z/2\to\Z$ byłoby
przekrojem $\pi$, którego
nie ma według uwagi po
Twierdzeniu~\ref{thm:ses-rozszczepienie}.
Zatem implikacja \emph{acykliczny
$\Rightarrow$ ściągalny} wymaga
dodatkowych założeń. Implikacja
odwrotna z
Twierdzenia~\ref{thm:homotopia-homologia}
działa nad każdym pierścieniem.
\end{przyklad}

\section{Ilorazy, ciągi względne i filtracje}
\label{sec:iloraz-filtracja}

Podkompleks $A_\bullet\subseteq C_\bullet$
oznacza podmoduły $A_n\subseteq C_n$
we wszystkich stopniach, takie że
$\partial_C(A_n)\subseteq A_{n-1}$.
Definiujemy \emph{kompleks ilorazowy}
$(C/A)_n=C_n/A_n$ i
$\partial_{C/A}[c]=[\partial_Cc]$.
Jeśli $[c]=[c+a]$ dla $a\in A_n$,
to $[\partial_Cc]=[\partial_Cc+\partial_Ca]$;
różniczka ilorazowa jest dobrze
określona i jej kwadrat wynosi zero.

\begin{stwierdzenie}[Długi ciąg ilorazu]
\label{prop:ses-iloraz}
Dla podkompleksu $A_\bullet\subseteq C_\bullet$
ciąg
\[
 0\longrightarrow A_\bullet
 \longrightarrow C_\bullet
 \longrightarrow C_\bullet/A_\bullet
 \longrightarrow0
\]
jest stopniowo dokładny. Dlatego istnieje
długi ciąg
\begin{equation}\label{eq:les-iloraz}
 \cdots\to H_n(A)\to H_n(C)\to H_n(C/A)
 \xrightarrow{\delta_n}H_{n-1}(A)\to\cdots.
\end{equation}
Jeśli $\xi\in H_n(C/A)$, to możemy
reprezentować ją elementem $c\in C_n$
z $\partial_Cc\in A_{n-1}$, a wtedy
$\delta_n\xi=[\partial_Cc]$. Tutaj $c$ najpierw wyznacza cykl
$c+A_n$ w kompleksie ilorazowym, a dopiero jego klasa modulo
brzegi jest elementem $\xi$; są to dwa różne przejścia do ilorazu.
\end{stwierdzenie}
\begin{proof}
W każdym stopniu włączenie $A_n\to C_n$
jest iniekcją, projekcja $C_n\to C_n/A_n$
jest surjekcją, a jej jądrem jest $A_n$.
Oba odwzorowania komutują z różniczkami
z definicji podkompleksu i ilorazu.
Twierdzenie~\ref{thm:les-homologia}
stosuje się więc bez dodatkowych
założeń. Warunek
$\partial_{C/A}[c]=[0]$ jest równoważny
$\partial_Cc\in A_{n-1}$; w tej sytuacji
konstrukcja~\eqref{eq:delta-wzor}
przy podniesieniu $c$ daje dokładnie
klasę $[\partial_Cc]$.
\end{proof}

Symbol $H_n(C,A):=H_n(C/A)$ będziemy
nazywać \emph{względną homologią
algebraiczną} kompleksu $C$ względem
podkompleksu $A$. Jest to definicja
czysto algebraiczna. W przyszłych
konstrukcjach topologicznych trzeba
osobno sprawdzić, że odpowiednie
łańcuchy tworzą podkompleks.

\begin{stwierdzenie}[Względna kohomologia nad ciałem]
\label{prop:wzgledna-kohomologia}
Jeśli $A_\bullet\subseteq C_\bullet$
jest podkompleksem przestrzeni liniowych
nad ciałem $\mathbb K$, to wzięcie
funkcjonałów daje krótki ciąg dokładny
kompleksów kołańcuchowych
\[
 0\to\operatorname{Hom}_{\mathbb K}
      (C/A,\mathbb K)
 \to\operatorname{Hom}_{\mathbb K}(C,\mathbb K)
 \to\operatorname{Hom}_{\mathbb K}(A,\mathbb K)
 \to0.
\]
Jeśli $H^n(C,A;\mathbb K)$ oznacza
kohomologię pierwszego z nich, to mamy
długi ciąg
\[
 \cdots\to H^n(C,A;\mathbb K)\to H^n(C;\mathbb K)
 \to H^n(A;\mathbb K)\to H^{n+1}(C,A;\mathbb K)
 \to\cdots .
\]
W tym zapisie $H^n(C;\mathbb K)$
oznacza kohomologię kompleksu
$\operatorname{Hom}_{\mathbb K}(C,\mathbb K)$,
a analogicznie dla $A$.
\end{stwierdzenie}
\begin{proof}
W stopniu $n$ pierwsze odwzorowanie składa
funkcjonał $C_n/A_n\to\mathbb K$
z projekcją $C_n\to C_n/A_n$,
więc jest iniektywne.
Drugie ogranicza funkcjonał na $C_n$
do podprzestrzeni $A_n$.
Jego jądrem są dokładnie funkcjonały
znikające na $A_n$, czyli faktoryzujące
się przez $C_n/A_n$; to obraz
pierwszego odwzorowania. Każdy funkcjonał
na $A_n$ przedłuża się do $C_n$:
wybieramy dopełnienie $A_n$
i stawiamy na nim wartość $0$.
Drugie odwzorowanie jest zatem surjektywne.
Odwzorowania komutują z różniczkami
kołańcuchowymi, bo są zdefiniowane
przez składanie z odwzorowaniami łańcuchowymi.
Stosujemy
Twierdzenie~\ref{thm:les-kohomologia}.
\end{proof}

\begin{stwierdzenie}[Trzy zagnieżdżone kompleksy]
\label{prop:les-potrojny}
Jeśli $A_\bullet\subseteq B_\bullet
\subseteq C_\bullet$, to istnieje
krótki ciąg dokładny kompleksów
\[
 0\longrightarrow B_\bullet/A_\bullet
 \longrightarrow C_\bullet/A_\bullet
 \longrightarrow C_\bullet/B_\bullet
 \longrightarrow0
\]
i wynikający z niego długi ciąg
\[
 \cdots\to H_n(B,A)\to H_n(C,A)\to H_n(C,B)
 \to H_{n-1}(B,A)\to\cdots.
\]
\end{stwierdzenie}
\begin{proof}
W stopniu $n$ pierwsze odwzorowanie
$[b]_A\mapsto[b]_A$ jest iniektywne:
jeśli $[b]_A=0$ w $C_n/A_n$, to
$b\in A_n$, więc była już zerowa w
$B_n/A_n$. Ostatnie odwzorowanie
$[c]_A\mapsto[c]_B$ jest dobrze
określone, bo $A_n\subseteq B_n$,
i surjektywne, bo każda klasa modulo
$B_n$ ma reprezentanta $c\in C_n$.
Jej jądro składa się z klas
$[c]_A$ z $c\in B_n$, czyli jest
obrazem pierwszego odwzorowania.
Różniczki zachowują wszystkie trzy
podkompleksy, więc są to odwzorowania
kompleksów. Stosujemy
Twierdzenie~\ref{thm:les-homologia}.
\end{proof}

\begin{definicja}[Filtracja kompleksu]
\index{filtracja kompleksu@Filtracja kompleksu}
\label{def:filtracja-kompleksu}
Przez \emph{skończoną filtrację rosnącą} rozumiemy tutaj ciąg
podkompleksów
\[
 0=F_{-1}C_\bullet\subseteq
 F_0C_\bullet\subseteq\cdots\subseteq
 F_NC_\bullet=C_\bullet.
\]
Iloraz $F_pC/F_{p-1}C$ opisuje
algebraicznie \emph{nową warstwę}
w kroku $p$. W każdym kroku mamy
krótki ciąg dokładny
$0\to F_{p-1}C\to F_pC\to
F_pC/F_{p-1}C\to0$
i odpowiadający mu długi ciąg homologii.
\end{definicja}

\begin{wniosek}[Zmiana skoncentrowana w jednym stopniu]
\label{wn:filtracja-jeden-stopien}
Niech $A\subseteq B$ będzie podkompleksem
w nieujemnych stopniach i niech
$H_j(B/A)=0$ dla $j\ne q$, $q\geq1$.
Wówczas $H_j(A)\to H_j(B)$ jest
izomorfizmem dla $j\notin\{q,q-1\}$,
a pozostałe zmiany opisuje ciąg
\begin{equation}\label{eq:jedna-warstwa}
 0\to H_q(A)\to H_q(B)\to H_q(B/A)
 \xrightarrow{\delta_q}H_{q-1}(A)
 \to H_{q-1}(B)\to0.
\end{equation}
\end{wniosek}
\begin{proof}
Dla dowolnego $j$ wytnijmy z~\eqref{eq:les-iloraz}
fragment
\[
 H_{j+1}(B/A)\to H_j(A)\to H_j(B)\to H_j(B/A).
\]
Jeśli oba zewnętrzne moduły są zerowe,
środkowe odwzorowanie jest jednocześnie
iniektywne i surjektywne.
Tak jest, gdy $j\ne q,q-1$.
Dla $j=q$ i $j=q-1$ zachowujemy
pełen fragment długiego ciągu.
Jego lewa skrajna strzałka wychodzi
z $H_{q+1}(B/A)=0$, a prawa
wchodzi do $H_{q-1}(B/A)=0$,
co daje dokładnie~\eqref{eq:jedna-warstwa}.
\end{proof}

\begin{uwaga}[Co mówi odwzorowanie łączące?]
W~\eqref{eq:jedna-warstwa} element
nowej warstwy może zwiększyć
$H_q(B)$, jeśli ma zerowy obraz
przez $\delta_q$. Może też sprawić,
że wcześniejsza klasa w
$H_{q-1}(A)$ stanie się brzegiem:
tak dzieje się, jeśli leży w
$\operatorname{im}\delta_q$.
Ten jeden fragment ciągu odróżnia
\emph{powstanie klasy} od
\emph{zabicia klasy}, bez odwoływania
się do szczególnej realizacji
geometrycznej kompleksu.
\end{uwaga}
\section{Suma podkompleksów i ciąg Mayera--Vietorisa}
\label{sec:mv-algebraiczny}

Przy rozbiciu obiektu na dwie części trzeba
uwzględnić elementy wspólne: bez tego
zliczylibyśmy je dwa razy. Poniższy wynik
jest czysto algebraicznym źródłem
późniejszych ciągów Mayera--Vietorisa.
Nie wymaga jeszcze żadnej konstrukcji
łańcuchów z przestrzeni topologicznej.

\begin{twierdzenie}[Algebraiczny Mayer--Vietoris]
\label{thm:mv-algebraiczny}
Niech $A_\bullet$ i $B_\bullet$ będą
podkompleksami $C_\bullet$. Wtedy
$A\cap B$ i $A+B$ też są podkompleksami,
a ciąg
\begin{equation}\label{eq:mv-ses}
 0\longrightarrow A\cap B
 \xrightarrow{\ j\ }A\oplus B
 \xrightarrow{\ q\ }A+B
 \longrightarrow0,
 \quad j(x)=(x,-x),\quad q(a,b)=a+b
\end{equation}
jest krótkim ciągiem dokładnym kompleksów.
W szczególności mamy długi ciąg dokładny
\begin{equation}\label{eq:mv-les}
 \cdots\to H_n(A\cap B)\to H_n(A)\oplus H_n(B)
 \to H_n(A+B)\xrightarrow{\delta}
 H_{n-1}(A\cap B)\to\cdots.
\end{equation}
Jeśli cykl $c\in(A+B)_n$ zapiszemy
$c=a+b$ z $a\in A_n$, $b\in B_n$,
to $\delta[c]=[\partial a]=[-\partial b]$.
\end{twierdzenie}
\begin{proof}
Jeśli $x\in A_n\cap B_n$, to
$\partial x\in A_{n-1}$ i $\partial x\in B_{n-1}$,
więc $\partial$ zachowuje przecięcie.
Podobnie $\partial(a+b)=\partial a+\partial b$ należy do
$A_{n-1}+B_{n-1}$, zatem suma
też jest podkompleksem.
Różniczka w $A\oplus B$ działa
składowo, dlatego
$\partial\,j(x)=(\partial x,-\partial x)=j(\partial x)$ i
$\partial\,q(a,b)=\partial(a+b)=q(\partial a,\partial b)$.
Są to odwzorowania łańcuchowe.

W każdym stopniu odwzorowanie $j$ jest
iniektywne, bo $j(x)=(0,0)$
implikuje $x=0$. Odwzorowanie $q$ jest
surjektywne na $A_n+B_n$ z samej
definicji sumy. Jeśli $q(a,b)=0$,
to $a=-b$ należy jednocześnie do
$A_n$ i do $B_n$, a para ma postać
$(a,-a)=j(a)$. Stąd
$\ker q=\operatorname{im}j$.
Twierdzenie~\ref{thm:les-homologia}
daje~\eqref{eq:mv-les}. Używamy przy tym kanonicznego izomorfizmu
$H_n(A\oplus B)\cong H_n(A)\oplus H_n(B)$:
różniczka działa składowo, więc cykle i brzegi są odpowiednio
$Z_n(A)\oplus Z_n(B)$ i $B_n(A)\oplus B_n(B)$,
a $[(a,b)]\mapsto([a],[b])$ jest izomorfizmem ich ilorazów.

Na koniec, jeśli $\partial(a+b)=0$, to
$\partial a=-\partial b$. Lewa strona należy do
$A_{n-1}$, prawa do $B_{n-1}$,
więc obie są elementem $A\cap B$.
W konstrukcji odwzorowania łączącego
podniesienie cyklu $c$ jest parą $(a,b)$,
a jej brzeg $(\partial a,\partial b)=(\partial a,-\partial a)$
jest obrazem $j(\partial a)$. To daje
podany wzór i objaśnia znak.
\end{proof}

\begin{przyklad}[Niezerowy brzeg z przecięcia]
\label{prz:mv-niezerowa-delta}
Nad ciałem $\mathbb K$ weźmy kompleks
z $C_1=\mathbb K u\oplus\mathbb K v$,
$C_0=\mathbb K x\oplus\mathbb K y$,
$\partial(u)=x-y$ i $\partial(v)=y-x$,
zerowy w pozostałych stopniach.
Niech $A_1=\mathbb K u$, $B_1=\mathbb K v$,
ale $A_0=B_0=C_0$. Są to podkompleksy,
$A+B=C$, zaś $A\cap B$ jest
skoncentrowany w stopniu $0$
i ma tam moduł $\mathbb K x\oplus\mathbb K y$.
Element $u+v$ jest cyklem w $C_1$.
Jego podniesienie do $A_1\oplus B_1$ to $(u,v)$,
którego brzeg wynosi
$(x-y,y-x)=j(x-y)$. Stąd
\[
 \delta[u+v]=[x-y]
 \in H_0(A\cap B)=\mathbb K x\oplus\mathbb K y,
\]
klasa niezerowa. Widzimy, jak brzeg
jednej części zostaje zniesiony
przez brzeg drugiej i powstaje
cykl w ich sumie.
\end{przyklad}

\section{Współczynniki, torsja i charakterystyka Eulera}
\label{sec:wspolczynniki-euler}

Wybór pierścienia $R$ wpływa na
wynik homologiczny. Kompleks
wolnych grup abelowych
można sprowadzić modulo $m$,
lecz homologia po tej operacji
nie polega zawsze wyłącznie
na redukcji dawnej homologii:
dochodzi torsja z poprzedniego stopnia.

\begin{twierdzenie}[Redukcja współczynników modulo $m$]
\label{thm:mod-m-homologia}
Niech $C_\bullet$ będzie kompleksem
wolnych grup abelowych, a $m\geq2$.
Różniczki wyznaczają kompleks
$(C/mC)_n=C_n/mC_n$. Dla każdego
$n$ istnieje krótki ciąg dokładny
\begin{equation}\label{eq:mod-m-homologia}
 0\longrightarrow H_n(C)/mH_n(C)
 \longrightarrow H_n(C/mC)
 \longrightarrow H_{n-1}(C)[m]
 \longrightarrow0,
\end{equation}
gdzie $H_{n-1}(C)[m]:=
\{x\in H_{n-1}(C):mx=0\}$
to $m$-torsja z Definicji~\ref{def:torsja}.
Nie twierdzimy, że ten ciąg ma
kanoniczne rozszczepienie.
\end{twierdzenie}
\begin{proof}
Wolna grupa abelowa nie ma
niezerowych elementów zabijanych
przez $m$. W każdym stopniu mamy
zatem dokładny ciąg
\[
 0\to C_n\xrightarrow{\ \times m\ }C_n
 \xrightarrow{\ \pi\ }C_n/mC_n\to0.
\]
Ponieważ $\partial(mx)=m\,\partial(x)$,
odwzorowania $\times m$ i $\pi$
komutują z różniczkami.
To krótki ciąg dokładny
kompleksów, więc
Twierdzenie~\ref{thm:les-homologia}
daje fragment
\[
 H_n(C)\xrightarrow{\ \times m\ }H_n(C)
 \xrightarrow{\ \pi_*\ }H_n(C/mC)
 \xrightarrow{\ \delta\ }H_{n-1}(C)
 \xrightarrow{\ \times m\ }H_{n-1}(C).
\]
Dokładność mówi
$\ker\pi_*=mH_n(C)$, a więc
$\pi_*$ wyznacza iniekcję
$H_n(C)/mH_n(C)\hookrightarrow H_n(C/mC)$
o obrazie $\ker\delta$.
Mówi też
$\operatorname{im}\delta
=\ker(\times m:H_{n-1}(C)\to H_{n-1}(C))
=H_{n-1}(C)[m]$.
Odwzorowanie $\delta$ po ograniczeniu
przeciwdziedziny do tego jądra
jest więc surjektywne.
To dokładnie trzy warunki
ciągu~\eqref{eq:mod-m-homologia}.
Warto zapisać ostatnią strzałkę jawnie. Jeśli $\bar c$ jest cyklem
modulo $m$ i $c\in C_n$ jest jego przedstawicielem, to
$\partial c=ma$ dla pewnego $a\in C_{n-1}$.
Element $a$ jest jednoznaczny, bo mnożenie przez $m$ jest iniektywne.
Ponadto $m\partial a=0$ wymusza $\partial a=0$,
a $m[a]=[\partial c]=0$. Zatem strzałka wysyła $[\bar c]$ na $[a]$.
Jest naturalna względem odwzorowań łańcuchowych, gdyż
$\partial f(c)=mf(a)$; pierwsza strzałka wysyła klasę cyklu
$[z]$ modulo $mH_n(C)$ na klasę $z$ modulo $m$.
\end{proof}

\begin{przyklad}[Nowa homologia po redukcji]
Dla $C_1=C_0=\Z$ i $\partial_1=\times2$
otrzymaliśmy $H_1(C)=0$ oraz
$H_0(C)=\Z/2$. Po redukcji modulo
$2$ różniczka $\times2$ staje się
zerowa, a więc
$H_1(C/2C)\cong H_0(C/2C)\cong\Z/2$.
Nowa klasa stopnia $1$ pochodzi
w~\eqref{eq:mod-m-homologia} z
$H_0(C)[2]=\Z/2$, choć wcześniejsze
$H_1(C)$ było zerowe.
\end{przyklad}

\begin{uwaga}[Iloczyn tensorowy modułów i rozszerzenie skalarów]
\label{uwaga:tensor-modulow}
Konstrukcję $M\otimes_R N$, rozszerzenie skalarów
i wpływ tensorowania na surjekcje i iniekcje
udowodniliśmy w Sekcji~\ref{sec:tensor-modulow-algebra}.
Tutaj stosujemy ją do grup łańcuchów; w szczególności
możemy pisać $C_n\otimes_{\Z}\mathbb F$ dla ciała
$\mathbb F$ zawierającego $\Q$.
\end{uwaga}

\begin{definicja}[Homologia ze współczynnikami w $\Q$, $\R$ lub $\C$]
\label{def:homologia-ciala-char-zero}
Niech $(C_\bullet,\partial)$ będzie kompleksem
wolnych grup abelowych, na przykład
kompleksem łańcuchów symplicjalnych.
Dla $\mathbb F\in\{\Q,\R,\C\}$ tworzymy
\[
 C_n(C;\mathbb F):=C_n\otimes_{\Z}\mathbb F,
 \qquad
 \partial_{\mathbb F}(c\otimes a)
       :=(\partial c)\otimes a.
\]
Jeśli $C_n$ ma bazę sympleksów, nowy kompleks
ma tę samą bazę, lecz można przy nich stawiać
współczynniki z $\mathbb F$ zamiast liczb całkowitych.
Jego homologię oznaczamy $H_n(C;\mathbb F)$;
dla kompleksu przypisanego przestrzeni $X$
piszemy $H_n(X;\mathbb F)$.
Piszemy też $H_n(C;\Z)=H_n(C)$ dla pierwotnego kompleksu.
Różniczka jest dobrze określona przez własność uniwersalną,
a $\partial_{\mathbb F}^2=0$ sprawdzamy na tensorach prostych.
\end{definicja}

\begin{twierdzenie}[Co zmienia przejście do charakterystyki zero]
\label{thm:homologia-char-zero}
Dla każdego kompleksu wolnych grup abelowych
i $\mathbb F=\Q,\R$ lub $\C$ istnieje naturalny
izomorfizm
\begin{equation}\label{eq:homologia-char-zero}
 H_n(C;\mathbb F)
 \cong H_n(C;\Z)\otimes_{\Z}\mathbb F.
\end{equation}
Jeśli $H_n(C;\Z)\cong\Z^{b_n}\oplus T_n$,
gdzie $T_n$ jest skończoną grupą torsyjną, to
$H_n(C;\mathbb F)\cong\mathbb F^{b_n}$.
W szczególności w tych trzech ciałach
pozostaje ta sama liczba $b_n$ niezależnych
klas, a torsja znika.
\end{twierdzenie}
\begin{proof}
Najpierw rozważmy $\mathbb F=\Q$.
Element $C_n\otimes_{\Z}\Q$ jest skończoną
sumą prostych tensorów; sprowadzając
mianowniki do wspólnego, zapisujemy go jako
$c/d$ z $c\in C_n$ i niezerowym
$d\in\Z$. Ponieważ każda grupa $C_j$
jest wolna, odwzorowania
$C_j\to C_j\otimes_{\Z}\Q$ są iniektywne.
Jeśli $c/d$ jest cyklem, to
$\partial c/d=0$, a więc $\partial c=0$.
Każda klasa racjonalna pochodzi zatem
całkowitego cyklu po dopuszczeniu mianownika:
odwzorowanie
\[
 H_n(C;\Z)\otimes_{\Z}\Q\longrightarrow H_n(C;\Q),
 \qquad [c]\otimes q\longmapsto[c\otimes q]
\]
jest surjektywne. Jest ono dobrze określone,
bo cykl pozostaje cyklem, brzeg pozostaje
brzegiem, a reguła jest dwuliniowa nad $\Z$.
Każdy tensor w dziedzinie ma postać $[c]\otimes(1/d)$
po sprowadzeniu skończenie wielu ułamków do wspólnego mianownika.

Sprawdźmy, czy nie utożsamia zbyt wielu klas.
Załóżmy, że cykl $c/d$ jest brzegiem
w kompleksie racjonalnym. Istnieją
$b\in C_{n+1}$ i $e\ne0$ takie, że
$c/d=\partial b/e$. Po pomnożeniu przez
$de$ i użyciu iniektywności
$C_n\to C_n\otimes\Q$ otrzymujemy
$ec=d\,\partial b=\partial(db)$
już w $C_n$. To znaczy, że $e[c]=0$
w $H_n(C;\Z)$. Po przejściu do $\Q$
liczba $e$ jest odwracalna, więc
$[c]\otimes(1/d)=0$. Odwzorowanie jest
również iniektywne i wzór został
udowodniony dla $\Q$.

Teraz $\mathbb F=\R$ albo $\C$.
Są to przestrzenie liniowe nad $\Q$, a
rozszerzenie skalarów daje izomorfizm
kompleksów
\[
 C_\bullet\otimes_{\Z}\mathbb F
 \cong (C_\bullet\otimes_{\Z}\Q)
             \otimes_{\Q}\mathbb F.
\]
Jawnie, $(c\otimes q)\otimes a\mapsto c\otimes qa$,
a odwrotnie $c\otimes a\mapsto(c\otimes1)\otimes a$.
Relacje tensorowe pokazują, że oba wzory są dobrze określone
i wzajemnie odwrotne; komutują też z różniczkami.
Wyjaśnijmy, dlaczego tensorowanie kompleksu
przestrzeni liniowych nad ciałem zachowuje
homologię. Oznaczmy na chwilę
$D_\bullet=C_\bullet\otimes_{\Z}\Q$,
$Z_n^\Q=\ker(\partial:D_n\to D_{n-1})$
oraz $B_n^\Q=\operatorname{im}(\partial:D_{n+1}\to D_n)$.
W każdym stopniu uzupełniamy bazy, aby zapisać
$Z_n^\Q=B_n^\Q\oplus V_n$ i
$D_n=Z_n^\Q\oplus L_n$.
Wtedy $\partial:L_n\to B_{n-1}^\Q$ jest
izomorfizmem: surjektywność wynika z definicji
$B_{n-1}^\Q$, a jądro to
$L_n\cap Z_n^\Q=0$. Ponadto
$V_n\cong Z_n^\Q/B_n^\Q=H_n(D)$.
Po tensorowaniu nad $\Q$ rozkłady nadal
są sumami prostymi, a ograniczenie
$\partial\otimes1:L_n\otimes\mathbb F
\to B_{n-1}^\Q\otimes\mathbb F$ pozostaje
izomorfizmem. Nowa homologia to więc
$V_n\otimes\mathbb F
\cong H_n(D)\otimes_{\Q}\mathbb F$.
Łącząc to z przypadkiem $\Q$, otrzymujemy
\eqref{eq:homologia-char-zero} dla $\R,\C$.
Choć dopełnienia ułatwiły dowód bijektywności, sam izomorfizm
jest kanoniczny: we wszystkich trzech przypadkach jest zadany
przez $[z]\otimes a\mapsto[z\otimes a]$.
Równość $f(z)\otimes a=(f\otimes\id)(z\otimes a)$ dowodzi
naturalności względem odwzorowań łańcuchowych.

Wreszcie $\Z\otimes_{\Z}\mathbb F\cong\mathbb F$.
Jeśli $mx=0$ w $T_n$, to
$m(x\otimes a)=0$; ponieważ $m\ne0$
ma odwrotność w $\mathbb F$,
$x\otimes a=0$. Zatem $T_n\otimes\mathbb F=0$,
a $b_n$ generatorów części wolnej daje
$\mathbb F^{b_n}$.
\end{proof}

Warto porównać to z Twierdzeniem~\ref{thm:mod-m-homologia}.
W ciele $\Q$, $\R$ lub $\C$ każda niezerowa
liczba całkowita staje się odwracalna.
Przy współczynnikach modulo $p$ liczba $p$
staje się natomiast \emph{zerem}; torsja
może wtedy dać nowe klasy także w innym stopniu.
Przejście do $\R$ lub $\C$ po obliczeniu
homologii nad $\Q$ jest już tylko
rozszerzeniem skalarów:
\[
 H_n(C;\R)\cong H_n(C;\Q)\otimes_{\Q}\R,
 \qquad
 H_n(C;\C)\cong H_n(C;\R)\otimes_{\R}\C.
\]

\begin{twierdzenie}[Algebraiczna charakterystyka Eulera]
\label{thm:euler-homologia}
Niech $C_\bullet$ będzie kompleksem
skończenie wymiarowych przestrzeni
liniowych nad ciałem $\mathbb K$
i niech tylko skończenie wiele $C_n$
będzie niezerowych. Wtedy
\begin{equation}\label{eq:euler-homologia}
 \sum_n(-1)^n\dim_{\mathbb K}C_n
 =\sum_n(-1)^n\dim_{\mathbb K}H_n(C).
\end{equation}
Lewą stronę nazywamy \emph{charakterystyką
Eulera kompleksu}.
\end{twierdzenie}
\begin{proof}
Odwzorowanie $\partial_n:C_n\to B_{n-1}$ jest
surjektywne z definicji obrazu,
a jego jądro to $Z_n$.
Wybór bazy $Z_n$ i jej uzupełnienie
do bazy $C_n$ pokazują
\[
 \dim C_n=\dim Z_n+\dim B_{n-1}.
\]
Podobnie projekcja $Z_n\to H_n=Z_n/B_n$
ma jądro $B_n$, więc
$\dim Z_n=\dim B_n+\dim H_n$.
Łącząc te równości, otrzymujemy
\[
 \dim C_n=\dim B_n+\dim H_n+\dim B_{n-1}.
\]
Mnożymy przez $(-1)^n$ i sumujemy
po $n$. Suma
$\sum_n(-1)^n\dim B_n$
znosi się z
$\sum_n(-1)^n\dim B_{n-1}$:
po podstawieniu $j=n-1$
druga jest równa
$-\sum_j(-1)^j\dim B_j$.
Skończona liczba niezerowych
wyrazów pozwala to wykonać bez
kłopotów ze zbieżnością.
Pozostaje~\eqref{eq:euler-homologia}.
\end{proof}

\begin{przyklad}[Wykrywanie brakującego wymiaru]
Niech nad $\mathbb K$ mamy tylko
$C_2\cong\mathbb K^2$, $C_1\cong\mathbb K^5$
i $C_0\cong\mathbb K^4$.
Niezależnie od szczegółów różniczek
spełniających $\partial^2=0$, musi zachodzić
\[
 \dim H_0-\dim H_1+\dim H_2
 =4-5+2=1.
\]
Nie można zatem uzyskać
acyklicznego kompleksu z takimi
wymiarami. Charakterystyka Eulera
nie mówi, w którym stopniu leży
niezerowa klasa; daje jednak
szybki test każdego obliczenia.
\end{przyklad}

\section{Porównywanie długich ciągów}
\label{sec:lemat-pieciu}

Czasem umiemy opisać cztery składniki
dwóch długich ciągów, a piąty jest trudny.
Następny lemat zamienia dane o odwzorowaniach
sąsiednich w informację o odwzorowaniu środkowym.
Przyda się do porównywania kompleksów
i różnych modeli tej samej homologii.

\begin{lemat}[Lemat pięciu]
\label{lem:pieciu}
Załóżmy, że mamy przemienny diagram
homomorfizmów $R$-modułów o dokładnych
wierszach
\[
\begin{array}{ccccccccc}
A_1&\longrightarrow&A_2&\longrightarrow&A_3&
 \longrightarrow&A_4&\longrightarrow&A_5\\
\downarrow f_1&&\downarrow f_2&&\downarrow f_3&&
\downarrow f_4&&\downarrow f_5\\
B_1&\longrightarrow&B_2&\longrightarrow&B_3&
 \longrightarrow&B_4&\longrightarrow&B_5 .
\end{array}
\]
Jeśli $f_1,f_2,f_4,f_5$ są izomorfizmami,
to $f_3$ także jest izomorfizmem.
\end{lemat}
\begin{proof}
Oznaczmy kolejne poziome odwzorowania w górnym
wierszu przez $u_1,u_2,u_3,u_4$,
a w dolnym przez $v_1,v_2,v_3,v_4$.
Przemienność oznacza
$f_{j+1}u_j=v_jf_j$.

\emph{Iniektywność.} Niech $x\in A_3$
spełnia $f_3x=0$.
Wtedy
$f_4(u_3x)=v_3(f_3x)=0$;
iniektywność $f_4$ daje $u_3x=0$.
Z dokładności w $A_3$ mamy
$x=u_2y$ dla pewnego $y\in A_2$.
Ponieważ
$v_2f_2y=f_3u_2y=f_3x=0$,
dokładność w $B_2$ daje
$f_2y=v_1z$ dla pewnego $z\in B_1$.
Surjektywność $f_1$ pozwala napisać
$z=f_1w$ dla $w\in A_1$.
Stąd $f_2y=v_1f_1w=f_2u_1w$.
Iniektywność $f_2$ daje $y=u_1w$,
a więc $x=u_2u_1w=0$.

\emph{Surjektywność.} Weźmy $z\in B_3$.
Surjektywność $f_4$ daje $a_4\in A_4$
z $f_4a_4=v_3z$. Teraz
$f_5u_4a_4=v_4f_4a_4=v_4v_3z=0$.
Z iniektywności $f_5$ wynika
$u_4a_4=0$. Dokładność w $A_4$
daje $a_4=u_3x$ dla pewnego
$x\in A_3$. Różnica
$z-f_3x$ leży w $\ker v_3$,
bo $v_3(z-f_3x)=v_3z-f_4u_3x=0$.
Z dokładności w $B_3$ jest więc
równa $v_2y$ dla pewnego $y\in B_2$.
Surjektywność $f_2$ daje
$y=f_2w$ dla $w\in A_2$, a wtedy
$z=f_3x+v_2f_2w
=f_3(x+u_2w)$.
Otrzymaliśmy surjektywność $f_3$.
\end{proof}

\section*{Dalsza lektura}
\addcontentsline{toc}{section}{Dalsza lektura}
Algebraiczny długi ciąg homologii i kohomologii
oraz homotopie kompleksów są opisane w
\emph{The Stacks Project}, rozdział
\emph{Homological Algebra},
\url{https://stacks.math.columbia.edu/tag/010V}.
Zastosowania długich ciągów do geometrii
topologicznej rozwija Allen Hatcher,
\emph{Algebraic Topology}, rozdział 2,
\url{https://pi.math.cornell.edu/~hatcher/AT/AT.pdf}.
Definicje potrzebne w obecnym rozdziale
i dowody podstawowych twierdzeń podaliśmy
bez używania jeszcze modeli
symplicjalnych ani komórkowych.
\clearpage
